% Chapter 8 draft for textbook inclusion. Assumes a textbook preamble with amsmath, amssymb, array, and booktabs.
\chapter{Degrees of Freedom, Subspaces, and Identifiability}
\label{chap:degrees-of-freedom-subspaces-identifiability}

\section{The Big Picture}

The previous chapter showed that linear regression is an orthogonal projection.  The response vector \(\mathbf{y}\in\mathbb{R}^n\) is projected onto the column space of the design matrix \(\mathbf{X}\).  The fitted values live in one subspace, and the residuals live in the orthogonal complement.

This chapter uses that geometry to explain three ideas that appear repeatedly in regression:
\begin{enumerate}
  \item \textbf{Degrees of freedom}: how many independent dimensions remain after we estimate model parameters.
  \item \textbf{Subspaces}: where fitted values, residuals, coefficient information, and non-identifiable directions live.
  \item \textbf{Identifiability}: when regression coefficients are uniquely determined by the data, and when many coefficient vectors produce the same fitted values.
\end{enumerate}

The central message is:
\begin{quote}
  Degrees of freedom are leftover dimensions.  Identifiability fails when the design matrix has a nonzero null space.
\end{quote}

For a full-rank multiple linear regression model with \(p\) predictors and an intercept,
\begin{equation}
  \mathbf{X}\in\mathbb{R}^{n\times(p+1)},
  \qquad
  \operatorname{rank}(\mathbf{X})=p+1,
\end{equation}
and the residual degrees of freedom are
\begin{equation}
  \boxed{n-p-1.}
\end{equation}
This is why the denominator of the residual variance estimator and the denominator degrees of freedom in the overall regression \(F\)-test use \(n-p-1\).

\paragraph{Math Foundations Connection.}
This chapter uses vector notation, the ones vector, dot products, orthogonality, matrix dimensions, matrix transposes, linear independence, bases, matrix-vector multiplication, left-inverse intuition, and null-space/subspace ideas.  These are covered in the Math Foundations textbook, Chapters~2, 5, 6, 7, 10, and~11.

\section{A First Definition: Degrees of Freedom as Independent Movement}

In data analysis, degrees of freedom are often introduced as a formula to memorize.  For example:
\begin{equation}
  s^2 = \frac{1}{n-1}\sum_{i=1}^n (y_i-\bar{y})^2.
\end{equation}
The natural question is: why \(n-1\)?

The geometric answer is that \(n\) observations define a vector in \(\mathbb{R}^n\).  Before estimating anything, the vector \(\mathbf{y}\) has \(n\) independent coordinates:
\begin{equation}
  \mathbf{y}=\begin{bmatrix}y_1\\y_2\\\vdots\\y_n\end{bmatrix}\in\mathbb{R}^n.
\end{equation}
Once we estimate the mean \(\bar{y}\), the residual deviations must satisfy one linear constraint:
\begin{equation}
  \sum_{i=1}^n (y_i-\bar{y}) = 0.
\end{equation}
One independent constraint removes one dimension.  Therefore, the residual variation around the sample mean has \(n-1\) degrees of freedom.

\paragraph{Reusable intuition.}
Each independently estimated parameter usually consumes one dimension.  The leftover dimensions are available for residual variation.

\section{Decomposing a Data Vector Around the Sample Mean}

The simplest projection problem is the sample mean.  Let
\begin{equation}
  \mathbf{1}=\begin{bmatrix}1\\1\\\vdots\\1\end{bmatrix}\in\mathbb{R}^n.
\end{equation}
The sample mean vector is
\begin{equation}
  \bar{y}\mathbf{1}.
\end{equation}
The residual vector around the sample mean is
\begin{equation}
  \mathbf{r}=\mathbf{y}-\bar{y}\mathbf{1}.
\end{equation}
Therefore,
\begin{equation}
  \boxed{\mathbf{y}=\bar{y}\mathbf{1}+\left(\mathbf{y}-\bar{y}\mathbf{1}\right).}
  \label{eq:ch8-mean-decomposition}
\end{equation}
This is a decomposition of the data vector into two orthogonal pieces:
\begin{enumerate}
  \item \(\bar{y}\mathbf{1}\), the projection onto the one-dimensional subspace \(\operatorname{span}\{\mathbf{1}\}\), and
  \item \(\mathbf{r}=\mathbf{y}-\bar{y}\mathbf{1}\), the residual vector in the orthogonal complement.
\end{enumerate}

To verify orthogonality, compute
\begin{align}
  \mathbf{1}^T\mathbf{r}
  &=\mathbf{1}^T(\mathbf{y}-\bar{y}\mathbf{1}) \\
  &=\sum_{i=1}^n y_i-\bar{y}\sum_{i=1}^n 1 \\
  &=n\bar{y}-n\bar{y} \\
  &=0.
\end{align}
Thus,
\begin{equation}
  \mathbf{r}\in \operatorname{span}\{\mathbf{1}\}^{\perp}.
\end{equation}
Since \(\operatorname{span}\{\mathbf{1}\}\) has dimension \(1\), its orthogonal complement in \(\mathbb{R}^n\) has dimension
\begin{equation}
  n-1.
\end{equation}
This is the geometric reason for the \(n-1\) denominator in sample variance.

\CoreCompress{
\section{Worked Example: Mean Decomposition with Three Observations}

Suppose
\begin{equation}
  \mathbf{y}=\begin{bmatrix}2\\5\\8\end{bmatrix}.
\end{equation}
The sample mean is
\begin{equation}
  \bar{y}=\frac{2+5+8}{3}=5.
\end{equation}
Then
\begin{equation}
  \bar{y}\mathbf{1}=5\begin{bmatrix}1\\1\\1\end{bmatrix}
  =\begin{bmatrix}5\\5\\5\end{bmatrix},
\end{equation}
and
\begin{equation}
  \mathbf{r}=\mathbf{y}-\bar{y}\mathbf{1}
  =\begin{bmatrix}2\\5\\8\end{bmatrix}-\begin{bmatrix}5\\5\\5\end{bmatrix}
  =\begin{bmatrix}-3\\0\\3\end{bmatrix}.
\end{equation}
Check the residual constraint:
\begin{equation}
  \mathbf{1}^T\mathbf{r}=(-3)+0+3=0.
\end{equation}
The residual vector lives on the plane
\begin{equation}
  \left\{\mathbf{v}\in\mathbb{R}^3: \mathbf{1}^T\mathbf{v}=0\right\}.
\end{equation}
That plane has dimension \(2\), which equals \(n-1=3-1\).
}{
\section{Worked Example: Mean Decomposition with Three Observations}

For \(\mathbf{y}=(2,5,8)^T\), the sample mean is \(\bar{y}=5\).  Therefore
\begin{equation}
  \bar{y}\mathbf{1}=\begin{bmatrix}5\\5\\5\end{bmatrix},
  \qquad
  \mathbf{r}=\mathbf{y}-\bar{y}\mathbf{1}=\begin{bmatrix}-3\\0\\3\end{bmatrix}.
\end{equation}
The residuals sum to zero, so they satisfy one linear constraint.  With three observations, that leaves \(3-1=2\) residual degrees of freedom.
}
\CoreCompress{
\section{The Mean Projection Matrix}

The sample mean is itself a projection.  Define
\begin{equation}
  \mathbf{H}_0=\frac{1}{n}\mathbf{1}\mathbf{1}^T.
\end{equation}
Then
\begin{equation}
  \mathbf{H}_0\mathbf{y}
  =\frac{1}{n}\mathbf{1}\mathbf{1}^T\mathbf{y}
  =\frac{1}{n}\mathbf{1}\sum_{i=1}^n y_i
  =\bar{y}\mathbf{1}.
\end{equation}
The residual-maker matrix for the mean-only model is
\begin{equation}
  \mathbf{M}_0=\mathbf{I}-\mathbf{H}_0.
\end{equation}
Then
\begin{equation}
  \mathbf{M}_0\mathbf{y}=\mathbf{y}-\bar{y}\mathbf{1}.
\end{equation}

The matrix \(\mathbf{H}_0\) is symmetric and idempotent, so it is an orthogonal projection matrix.  Its rank and trace are both \(1\):
\begin{equation}
  \operatorname{rank}(\mathbf{H}_0)=\operatorname{tr}(\mathbf{H}_0)=1.
\end{equation}
Likewise,
\begin{equation}
  \operatorname{rank}(\mathbf{M}_0)=\operatorname{tr}(\mathbf{M}_0)=n-1.
\end{equation}
This trace perspective becomes useful in regression because the hat matrix for linear regression is also a symmetric idempotent projection matrix.
}{
\section{The Mean Projection Matrix}

The mean-only fit can also be written as a projection:
\begin{equation}
  \widehat{\mathbf{y}}_0=\bar{y}\mathbf{1}=\mathbf{H}_0\mathbf{y},
  \qquad
  \mathbf{H}_0=\frac{1}{n}\mathbf{1}\mathbf{1}^T.
\end{equation}
This is the one-dimensional version of the hat-matrix idea.  The full reference text develops the symmetry, idempotence, rank, and trace details.
}
\section{Regression Decomposes \texorpdfstring{\(\mathbb{R}^n\)}{R n} Into Two Orthogonal Pieces}

For multiple linear regression with an intercept and \(p\) predictors, write
\begin{equation}
  \mathbf{y}=\mathbf{X}\boldsymbol{\beta}+\boldsymbol{\epsilon},
  \qquad
  \mathbf{X}\in\mathbb{R}^{n\times(p+1)}.
\end{equation}
The fitted values are
\begin{equation}
  \widehat{\mathbf{y}}=\mathbf{H}\mathbf{y},
  \qquad
  \mathbf{H}=\mathbf{X}(\mathbf{X}^T\mathbf{X})^{-1}\mathbf{X}^T,
\end{equation}
when \(\mathbf{X}\) has full column rank.

The residuals are
\begin{equation}
  \mathbf{e}=\mathbf{y}-\widehat{\mathbf{y}}=(\mathbf{I}-\mathbf{H})\mathbf{y}=\mathbf{M}\mathbf{y}.
\end{equation}
Regression decomposes the observation space as
\begin{equation}
  \boxed{\mathbb{R}^n=\operatorname{Col}(\mathbf{X})\oplus \operatorname{Null}(\mathbf{X}^T).}
  \label{eq:ch8-column-leftnull-decomposition}
\end{equation}
Here \(\oplus\) means a direct sum: every vector in \(\mathbb{R}^n\) can be written uniquely as a piece in \(\operatorname{Col}(\mathbf{X})\) plus a piece in \(\operatorname{Null}(\mathbf{X}^T)\).  These two pieces are orthogonal.

In regression terms:
\begin{align}
  \widehat{\mathbf{y}} &\in \operatorname{Col}(\mathbf{X}), \\
  \mathbf{e} &\in \operatorname{Null}(\mathbf{X}^T), \\
  \widehat{\mathbf{y}}^T\mathbf{e} &=0.
\end{align}
The fitted values live in the column space.  The residuals live in the left null space.

\section{Residual Degrees of Freedom}

Assume \(\mathbf{X}\) has full column rank.  Then
\begin{equation}
  \operatorname{rank}(\mathbf{X})=p+1.
\end{equation}
Therefore,
\begin{equation}
  \dim \operatorname{Col}(\mathbf{X})=p+1.
\end{equation}
The column space consumes \(p+1\) dimensions of the original \(n\)-dimensional observation space.  The leftover dimensions are the residual degrees of freedom:
\begin{equation}
  \boxed{\operatorname{df}_{\text{resid}}=n-(p+1)=n-p-1.}
\end{equation}

This explains the residual variance estimator:
\begin{equation}
  \widehat{\sigma}^2
  =\frac{\text{RSS}}{n-p-1}
  =\frac{\mathbf{e}^T\mathbf{e}}{n-p-1}.
\end{equation}
It also explains the denominator degrees of freedom in the overall \(F\)-test for MLR.

\paragraph{Special cases.}
\begin{center}
\begin{tabular}{|p{0.31\textwidth}|p{0.28\textwidth}|p{0.24\textwidth}|}
\hline
\textbf{Model} & \textbf{Fitted subspace dimension} & \textbf{Residual df} \\
\hline
Mean-only model & \(1\) & \(n-1\) \\
\hline
SLR with intercept & \(2\) & \(n-2\) \\
\hline
MLR with intercept and \(p\) predictors & \(p+1\) & \(n-p-1\) \\
\hline
Full rank model with design rank \(r\) & \(r\) & \(n-r\) \\
\hline
\end{tabular}
\end{center}

\CoreCompress{
\section{A Subtle Point: Model Degrees of Freedom Versus Fitted-Subspace Dimension}

There are two closely related uses of degrees of freedom in regression:
\begin{enumerate}
  \item \textbf{Fitted-subspace dimension}: with an intercept and \(p\) predictors, \(\operatorname{Col}(\mathbf{X})\) has dimension \(p+1\).
  \item \textbf{Model improvement degrees of freedom}: when comparing the full model to the intercept-only model, the predictors add \(p\) dimensions beyond the mean direction.
\end{enumerate}

That is why an ANOVA-style regression decomposition often uses
\begin{equation}
  \text{Total df}=n-1,
  \qquad
  \text{Regression df}=p,
  \qquad
  \text{Residual df}=n-p-1.
\end{equation}
The \(n-1\) total degrees of freedom already subtract the mean direction.  The regression degrees of freedom \(p\) count the additional fitted directions beyond the intercept-only model.
}{
\section{A Subtle Point: Model Degrees of Freedom Versus Fitted-Subspace Dimension}

Two counts appear in regression.  The fitted subspace has dimension \(p+1\) when the model has an intercept and \(p\) predictors.  But the overall regression \(F\)-test compares the full model to the intercept-only model, so its numerator degrees of freedom are \(p\), not \(p+1\).  The residual degrees of freedom are \(n-p-1\).
}
\CoreCompress{
\section{Trace of the Hat Matrix}

For a full-rank linear regression model, the hat matrix is
\begin{equation}
  \mathbf{H}=\mathbf{X}(\mathbf{X}^T\mathbf{X})^{-1}\mathbf{X}^T.
\end{equation}
It is symmetric and idempotent.  For symmetric idempotent matrices, the trace equals the rank.  Therefore,
\begin{equation}
  \boxed{\operatorname{tr}(\mathbf{H})=\operatorname{rank}(\mathbf{H})=\operatorname{rank}(\mathbf{X}).}
\end{equation}
If \(\mathbf{X}\) has full column rank, this becomes
\begin{equation}
  \operatorname{tr}(\mathbf{H})=p+1.
\end{equation}
The residual-maker matrix is \(\mathbf{M}=\mathbf{I}-\mathbf{H}\), so
\begin{align}
  \operatorname{tr}(\mathbf{M})
  &=\operatorname{tr}(\mathbf{I})-\operatorname{tr}(\mathbf{H}) \\
  &=n-(p+1) \\
  &=n-p-1.
\end{align}
This is another way to see residual degrees of freedom.

\paragraph{Connection to leverage.}
The diagonal entries of \(\mathbf{H}\), denoted \(h_{ii}\), are leverages.  Since the trace is the sum of diagonal entries,
\begin{equation}
  \sum_{i=1}^n h_{ii}=p+1.
\end{equation}
The average leverage is therefore
\begin{equation}
  \bar{h}=\frac{p+1}{n}.
\end{equation}
Large leverage means an observation is using more than its average share of the fitted-subspace degrees of freedom.
}{
\section{Trace of the Hat Matrix}

For a full-rank design matrix, the hat matrix projects onto \(\operatorname{Col}(\mathbf{X})\).  Its trace equals the fitted-subspace dimension:
\begin{equation}
  \operatorname{tr}(\mathbf{H})=\operatorname{rank}(\mathbf{H})=p+1.
\end{equation}
Thus the average leverage is \((p+1)/n\), and the residual-maker matrix has trace \(n-p-1\), matching the residual degrees of freedom.
}
\section{The Big Four Subspaces of the Design Matrix}

Let
\begin{equation}
  \mathbf{X}\in\mathbb{R}^{n\times(p+1)}
\end{equation}
with rank \(r\).  Linear algebra associates four fundamental subspaces with \(\mathbf{X}\).  In regression, these subspaces have direct meanings.

\begin{center}
\begin{tabular}{|p{0.24\textwidth}|p{0.26\textwidth}|p{0.20\textwidth}|p{0.16\textwidth}|}
\hline
\textbf{Subspace} & \textbf{Where it lives} & \textbf{Regression meaning} & \textbf{Dimension} \\
\hline
\(\operatorname{Col}(\mathbf{X})\) & \(\mathbb{R}^n\) & Fitted values \(\widehat{\mathbf{y}}\) & \(r\) \\
\hline
\(\operatorname{Null}(\mathbf{X}^T)\) & \(\mathbb{R}^n\) & Residual vectors \(\mathbf{e}\) & \(n-r\) \\
\hline
\(\operatorname{Row}(\mathbf{X})=\operatorname{Col}(\mathbf{X}^T)\) & \(\mathbb{R}^{p+1}\) & Identifiable coefficient directions & \(r\) \\
\hline
\(\operatorname{Null}(\mathbf{X})\) & \(\mathbb{R}^{p+1}\) & Non-identifiable coefficient directions & \((p+1)-r\) \\
\hline
\end{tabular}
\end{center}

The first two subspaces live in observation space \(\mathbb{R}^n\).  The last two live in coefficient space \(\mathbb{R}^{p+1}\).  This distinction prevents a common mistake: residuals and fitted values are vectors with \(n\) entries, while coefficient directions have \(p+1\) entries.

\section{Column Space and Left Null Space}

The fitted values satisfy
\begin{equation}
  \widehat{\mathbf{y}}=\mathbf{X}\widehat{\boldsymbol{\beta}}.
\end{equation}
Therefore, \(\widehat{\mathbf{y}}\) is a linear combination of the columns of \(\mathbf{X}\):
\begin{equation}
  \widehat{\mathbf{y}}\in\operatorname{Col}(\mathbf{X}).
\end{equation}
The residual vector satisfies the normal equations:
\begin{equation}
  \mathbf{X}^T\mathbf{e}=\mathbf{0}.
\end{equation}
Therefore,
\begin{equation}
  \mathbf{e}\in\operatorname{Null}(\mathbf{X}^T).
\end{equation}
Since \(\operatorname{Null}(\mathbf{X}^T)\) is the orthogonal complement of \(\operatorname{Col}(\mathbf{X})\), every residual vector is orthogonal to every fitted-value direction.

This is the geometric core of regression diagnostics:
\begin{equation}
  \mathbf{y}=\widehat{\mathbf{y}}+\mathbf{e},
  \qquad
  \widehat{\mathbf{y}}\perp \mathbf{e}.
\end{equation}

\section{Row Space and Null Space}

The row space and null space of \(\mathbf{X}\) live in coefficient space.  They answer a different question:
\begin{quote}
  Which coefficient directions change fitted values, and which coefficient directions are invisible to the data?
\end{quote}

A vector \(\mathbf{v}\in\mathbb{R}^{p+1}\) is in the null space of \(\mathbf{X}\) if
\begin{equation}
  \mathbf{X}\mathbf{v}=\mathbf{0}.
\end{equation}
If \(\mathbf{v}\ne\mathbf{0}\), then changing the coefficient vector in the direction \(\mathbf{v}\) does not change fitted values:
\begin{equation}
  \mathbf{X}(\boldsymbol{\beta}+c\mathbf{v})
  =\mathbf{X}\boldsymbol{\beta}+c\mathbf{X}\mathbf{v}
  =\mathbf{X}\boldsymbol{\beta}
\end{equation}
for any scalar \(c\).  Therefore, \(\mathbf{v}\) is a non-identifiable direction.

The row space contains coefficient directions that the data can see.  The null space contains coefficient directions that the data cannot see.

\section{Identifiability}

A parameter is identifiable when different parameter values imply different observable behavior.  In linear regression, the coefficient vector \(\boldsymbol{\beta}\) is identifiable when the mapping
\begin{equation}
  \boldsymbol{\beta}\mapsto \mathbf{X}\boldsymbol{\beta}
\end{equation}
is one-to-one.

This happens exactly when
\begin{equation}
  \operatorname{Null}(\mathbf{X})=\{\mathbf{0}\}.
\end{equation}
Equivalently, the columns of \(\mathbf{X}\) must be linearly independent.

For \(\mathbf{X}\in\mathbb{R}^{n\times(p+1)}\), the following statements are equivalent:
\begin{enumerate}
  \item \(\boldsymbol{\beta}\) is identifiable.
  \item \(\operatorname{Null}(\mathbf{X})=\{\mathbf{0}\}\).
  \item The columns of \(\mathbf{X}\) are linearly independent.
  \item \(\operatorname{rank}(\mathbf{X})=p+1\).
  \item \(\mathbf{X}^T\mathbf{X}\) is invertible.
  \item The least squares coefficient vector
  \begin{equation}
    \widehat{\boldsymbol{\beta}}=(\mathbf{X}^T\mathbf{X})^{-1}\mathbf{X}^T\mathbf{y}
  \end{equation}
  is unique.
\end{enumerate}

\paragraph{Why \(\mathbf{X}^T\mathbf{X}\) matters.}
For any \(\mathbf{v}\),
\begin{equation}
  \mathbf{v}^T\mathbf{X}^T\mathbf{X}\mathbf{v}=\|\mathbf{X}\mathbf{v}\|^2.
\end{equation}
If \(\mathbf{X}\mathbf{v}=\mathbf{0}\) for some nonzero \(\mathbf{v}\), then \(\mathbf{X}^T\mathbf{X}\) cannot be invertible.  Conversely, if \(\mathbf{X}^T\mathbf{X}\) is singular, then there is a nonzero \(\mathbf{v}\) such that \(\mathbf{X}\mathbf{v}=\mathbf{0}\).  Thus, \(\mathbf{X}^T\mathbf{X}\) is invertible exactly when \(\mathbf{X}\) has no nonzero null directions.

\section{What Fails When the Model Is Not Identified?}

When \(\operatorname{Null}(\mathbf{X})\ne\{\mathbf{0}\}\), there are infinitely many coefficient vectors that produce the same fitted values.  If \(\mathbf{v}\in\operatorname{Null}(\mathbf{X})\), then
\begin{equation}
  \mathbf{X}\widehat{\boldsymbol{\beta}}
  =\mathbf{X}(\widehat{\boldsymbol{\beta}}+c\mathbf{v})
\end{equation}
for every scalar \(c\).  The data cannot distinguish \(\widehat{\boldsymbol{\beta}}\) from \(\widehat{\boldsymbol{\beta}}+c\mathbf{v}\).

Two important consequences follow:
\begin{enumerate}
  \item The coefficient vector is not unique.
  \item The fitted values may still be unique, because the projection of \(\mathbf{y}\) onto \(\operatorname{Col}(\mathbf{X})\) is unique.
\end{enumerate}

Thus, non-identifiability is primarily a coefficient problem.  The model may still make fitted predictions, but the individual coefficients cannot be interpreted uniquely.

\CoreCompress{
\section{Non-Identifiability Example: Celsius and Fahrenheit}

Suppose a model includes an intercept, temperature measured in Celsius, and temperature measured in Fahrenheit:
\begin{equation}
  y_i=\beta_0+\beta_C C_i+\beta_F F_i+\epsilon_i.
\end{equation}
The two temperature predictors are exactly related:
\begin{equation}
  F_i=\frac{9}{5}C_i+32.
\end{equation}
The design matrix is
\begin{equation}
  \mathbf{X}=\begin{bmatrix}
    1 & C_1 & F_1 \\
    1 & C_2 & F_2 \\
    \vdots & \vdots & \vdots \\
    1 & C_n & F_n
  \end{bmatrix}.
\end{equation}
The columns are linearly dependent because
\begin{equation}
  \mathbf{x}_F-\frac{9}{5}\mathbf{x}_C-32\mathbf{1}=\mathbf{0}.
\end{equation}
Equivalently,
\begin{equation}
  \mathbf{X}\begin{bmatrix}-32\\-9/5\\1\end{bmatrix}=\mathbf{0}.
\end{equation}
Therefore,
\begin{equation}
  \begin{bmatrix}-32\\-9/5\\1\end{bmatrix}\in\operatorname{Null}(\mathbf{X}).
\end{equation}
This means that coefficient vectors can move in this direction without changing predictions.

For example, the coefficient vector
\begin{equation}
  \boldsymbol{\beta}=\begin{bmatrix}10\\2\\0\end{bmatrix}
\end{equation}
gives
\begin{equation}
  \widehat{y}=10+2C.
\end{equation}
But adding the null-space vector gives
\begin{equation}
  \boldsymbol{\beta}^{\star}
  =\begin{bmatrix}10\\2\\0\end{bmatrix}
  +\begin{bmatrix}-32\\-9/5\\1\end{bmatrix}
  =\begin{bmatrix}-22\\1/5\\1\end{bmatrix}.
\end{equation}
This gives
\begin{align}
  \widehat{y}^{\star}
  &=-22+\frac{1}{5}C+F \\
  &=-22+\frac{1}{5}C+\left(\frac{9}{5}C+32\right) \\
  &=10+2C.
\end{align}
The predictions are identical, but the coefficient interpretations are different.  Therefore, the individual Celsius and Fahrenheit coefficients are not identified.
}{
\section{Non-Identifiability Example: Celsius and Fahrenheit}

If a model includes an intercept, Celsius, and Fahrenheit, then one column is an exact affine transformation of another:
\begin{equation}
  F_i = 1.8C_i+32.
\end{equation}
In column form,
\begin{equation}
  \mathbf{f}=1.8\mathbf{c}+32\mathbf{1}.
\end{equation}
Thus the columns \(\mathbf{1}\), \(\mathbf{c}\), and \(\mathbf{f}\) are linearly dependent.  The fitted values may still be unique, but the individual Celsius and Fahrenheit coefficients are not.
}
\section{Common Sources of Non-Identifiability}

Exact non-identifiability usually comes from exact linear dependence among columns of the design matrix.  Common examples include:
\begin{enumerate}
  \item \textbf{Duplicate variables}: including the same predictor twice.
  \item \textbf{Unit conversions}: including meters and feet, or Celsius and Fahrenheit, in a model with an intercept.
  \item \textbf{The dummy variable trap}: including an intercept and all category indicator columns.
  \item \textbf{Derived totals}: including all component variables and their exact sum.
  \item \textbf{Too many predictors}: if \(p+1>n\), then \(\mathbf{X}\) cannot have full column rank.
\end{enumerate}

\paragraph{Dummy variable trap.}
Suppose a categorical variable has three levels: Army, Navy, and Marine Corps.  Let \(\mathbf{d}_A,\mathbf{d}_N,\mathbf{d}_M\) be the three indicator columns.  For every row,
\begin{equation}
  d_{iA}+d_{iN}+d_{iM}=1.
\end{equation}
Thus,
\begin{equation}
  \mathbf{d}_A+\mathbf{d}_N+\mathbf{d}_M=\mathbf{1}.
\end{equation}
If the design matrix includes \(\mathbf{1}\), \(\mathbf{d}_A\), \(\mathbf{d}_N\), and \(\mathbf{d}_M\), then its columns are linearly dependent.  The usual repair is to include the intercept and drop one category indicator as the reference group.

\section{Near Non-Identifiability: Multicollinearity}

Exact linear dependence causes exact non-identifiability.  Near linear dependence causes instability.  This is multicollinearity.

If one column of \(\mathbf{X}\) is almost a linear combination of other columns, then \(\mathbf{X}^T\mathbf{X}\) may still be invertible, but it is close to singular.  The least squares coefficients exist, but small changes in the data can produce large changes in the coefficient estimates.

This matters because:
\begin{enumerate}
  \item coefficient standard errors become large,
  \item signs and magnitudes can be unstable,
  \item individual coefficient interpretations become fragile, and
  \item predictions may remain reasonable even while coefficients are hard to interpret.
\end{enumerate}

\paragraph{Operational distinction.}
\begin{center}
\begin{tabular}{|p{0.30\textwidth}|p{0.34\textwidth}|p{0.24\textwidth}|}
\hline
\textbf{Condition} & \textbf{Design-matrix geometry} & \textbf{Consequence} \\
\hline
Exact dependence & Nonzero \(\operatorname{Null}(\mathbf{X})\) & Coefficients not unique \\
\hline
Near dependence & Columns nearly dependent & Coefficients unstable \\
\hline
Full rank and well-conditioned & Columns independent and not nearly redundant & Coefficients more stable \\
\hline
\end{tabular}
\end{center}

\CoreCompress{
\section{Worked Example: Counting Regression Degrees of Freedom}

Suppose we have \(n=40\) observations and fit an MLR model with an intercept and \(p=5\) predictors.  If the design matrix has full column rank, then
\begin{equation}
  \mathbf{X}\in\mathbb{R}^{40\times 6},
  \qquad
  \operatorname{rank}(\mathbf{X})=6.
\end{equation}
The fitted values live in a \(6\)-dimensional subspace:
\begin{equation}
  \dim\operatorname{Col}(\mathbf{X})=6.
\end{equation}
The residuals live in the orthogonal complement:
\begin{equation}
  \dim\operatorname{Null}(\mathbf{X}^T)=40-6=34.
\end{equation}
Therefore,
\begin{equation}
  \boxed{\operatorname{df}_{\text{resid}}=34.}
\end{equation}

For the overall \(F\)-test comparing the full model to the intercept-only model,
\begin{equation}
  \text{numerator df}=5,
  \qquad
  \text{denominator df}=34.
\end{equation}
The numerator has \(5\) degrees of freedom because the full model adds five predictor directions beyond the mean-only model.
}{
\section{Worked Example: Counting Regression Degrees of Freedom}

If \(n=40\) and we fit an MLR model with an intercept and \(p=5\) predictors, then the full-rank design matrix has \(p+1=6\) columns.  The fitted values live in a 6-dimensional subspace, and the residuals have
\begin{equation}
  40-6=34
\end{equation}
residual degrees of freedom.  The overall \(F\)-test has numerator degrees of freedom \(5\) and denominator degrees of freedom \(34\).
}
\paragraph{Coding companion reference.}
Package-specific Python/R commands for checking rank, residual degrees of freedom, and non-identifiability have been moved to the separate Coding and Software Cookbook companion.  The main chapter keeps the degrees-of-freedom and subspace interpretation.

\section{Common Mistakes}

\begin{enumerate}
  \item \textbf{Forgetting the intercept.}  With an intercept and \(p\) predictors, the design matrix has \(p+1\) columns, not \(p\).
  \item \textbf{Mixing up fitted-space dimension and regression \(F\)-test numerator df.}  The fitted subspace has dimension \(p+1\), but the overall \(F\)-test numerator df is \(p\) when comparing to the intercept-only model.
  \item \textbf{Thinking residuals live in predictor space.}  Residuals are \(n\)-vectors.  They live in \(\mathbb{R}^n\), not in \(\mathbb{R}^p\).
  \item \textbf{Thinking coefficient null-space directions are residual directions.}  \(\operatorname{Null}(\mathbf{X})\) lives in coefficient space, while \(\operatorname{Null}(\mathbf{X}^T)\) lives in observation space.
  \item \textbf{Using \((\mathbf{X}^T\mathbf{X})^{-1}\) when \(\mathbf{X}\) is rank deficient.}  If columns are dependent, the inverse does not exist.
  \item \textbf{Assuming fitted values are not unique when coefficients are not unique.}  Coefficients may be non-unique while the projection \(\widehat{\mathbf{y}}\) is still unique.
  \item \textbf{Checking only pairwise correlations for identifiability.}  A column can be a combination of several others even if no single pair looks perfectly correlated.
  \item \textbf{Including every dummy variable plus an intercept.}  This creates exact linear dependence.
  \item \textbf{Treating multicollinearity as a prediction failure only.}  Multicollinearity often harms coefficient interpretation before it visibly harms prediction.
\end{enumerate}

\CoreCompress{
\section{Chapter Summary}

Degrees of freedom are best understood as dimensions.  An \(n\)-vector of observations begins with \(n\) dimensions.  Estimating the sample mean projects \(\mathbf{y}\) onto \(\operatorname{span}\{\mathbf{1}\}\), leaving residual variation in an \((n-1)\)-dimensional orthogonal complement.

Regression generalizes the same idea.  The design matrix \(\mathbf{X}\) defines the fitted-value subspace.  If \(\mathbf{X}\in\mathbb{R}^{n\times(p+1)}\) has full column rank, then fitted values live in a \((p+1)\)-dimensional column space, and residuals live in an \((n-p-1)\)-dimensional orthogonal complement.  Thus,
\begin{equation}
  \operatorname{df}_{\text{resid}}=n-p-1.
\end{equation}

The four fundamental subspaces of \(\mathbf{X}\) organize the regression geometry:
\begin{align}
  \operatorname{Col}(\mathbf{X}) &\quad \text{contains fitted values}, \\
  \operatorname{Null}(\mathbf{X}^T) &\quad \text{contains residuals}, \\
  \operatorname{Row}(\mathbf{X}) &\quad \text{contains identifiable coefficient directions}, \\
  \operatorname{Null}(\mathbf{X}) &\quad \text{contains non-identifiable coefficient directions}.
\end{align}

Identifiability requires \(\operatorname{Null}(\mathbf{X})=\{\mathbf{0}\}\).  If the design matrix has linearly dependent columns, then infinitely many coefficient vectors can produce the same fitted values.  The Celsius/Fahrenheit example shows this directly: including an intercept, Celsius, and Fahrenheit creates exact linear dependence because Fahrenheit is an affine transformation of Celsius.

\section{Practice and Workflow}
\subsection*{Focused Study Checklist}

Use this as a final check after a degrees-of-freedom or identifiability problem.

\begin{enumerate}
  \item Can you identify \(n\), \(p\), and the dimensions of \(\mathbf{X}\)?
  \item Can you determine \(r=\operatorname{rank}(\mathbf{X})\)?
  \item Can you state fitted-space degrees of freedom \(r\) and residual degrees of freedom \(n-r\)?
  \item In the full-rank intercept model, can you simplify residual degrees of freedom to \(n-p-1\)?
  \item Can you explain why identifiability requires \(\operatorname{Null}(\mathbf{X})=\{\mathbf{0}\}\)?
  \item Can you recognize exact dependence, such as redundant columns or the dummy-variable trap?
  \item Can you distinguish exact non-identifiability from near dependence or multicollinearity?
\end{enumerate}
}{
\section{Chapter Summary}

Degrees of freedom are leftover dimensions.  An \(n\)-vector of observations begins with \(n\) independent coordinates.  Estimating the mean consumes one dimension, leaving \(n-1\) dimensions for residual variation.

Regression generalizes the same idea.  For a full-rank design matrix with an intercept and \(p\) predictors,
\begin{equation}
  \operatorname{df}_{\text{resid}}=n-p-1.
\end{equation}
The fitted values live in \(\operatorname{Col}(\mathbf{X})\), and the residuals live in \(\operatorname{Null}(\mathbf{X}^T)\).

Identifiability requires \(\operatorname{Null}(\mathbf{X})=\{\mathbf{0}\}\).  If the design matrix has dependent columns, multiple coefficient vectors can produce the same fitted values.  Exact dependence creates non-identifiability; near dependence creates unstable coefficient estimates.

\paragraph{Can you do this?}
You should be able to count residual degrees of freedom, distinguish fitted-space dimension from regression-test numerator degrees of freedom, identify the four subspaces in regression language, and recognize common non-identifiability patterns such as duplicated variables, unit conversions, and the dummy-variable trap.  The full reference text contains the reusable study template; the consolidated Math Foundations source map appears in Appendix~\ref{app:math-foundations-source-map}.
}
